\section{Shifted \texorpdfstring{$\cD$}{D}-modules and shifted deformation quantization}
\label{sec:dmod}

The theory of $\cD$-modules, modules over the sheaf of differential operators, is closely
related to the theory of quasi-coherent sheaves in Section \ref{sec:QC}. Of its many treatments
\cite{HTT, GR:v2}, we use the one in which $\cD$-modules on $X$ are ind-coherent sheaves on the
de Rham stack $X_\dR$ \cite{GR:v2}.
This was further generalized to the notion of a \emph{shifted} differential operator (and
correspondingly shifted $\cD$-module) in
\cite{Ber:Sph,Ber:center}, using the sheaf theory $\IndCoh_0$ of Arinkin--Gaitsgory
\cite{AG:crystal}.
Ho and Li \cite{HoLi} generalized this to the $\EE_n$-Hecke categories
$\mathrm{Hecke}_n\left(Y,Z\right)$ attached to a map of stacks $Y\to Z$, which recover
shifted $\cD$-modules for $Z=\pt$, and computed their factorization homology over
arbitrary manifolds.
We will work in this setting.
Using their excision theorem, we compute the higher centers of $\DMod^2\left(X\right)$ as
factorization homology over spheres (Theorem \ref{thm:centers}), alongside the values of
the associated TQFT on closed surfaces \cite{Ber:Sph,HoLi}.
The upshot (Theorem \ref{thm:BW}) is that for $X = V\left[1\right]$, with $V$ a non-zero vector
space, the $\EE_2$-category $\DMod^2\left(X\right)$ is $3$-dualizable and \emph{non-degenerate}
(Definition \ref{defn:inv_E2}) in the Morita $\left(\infty , 4\right)$-category of
$\EE_2$-monoidal categories (see Section \ref{sec:dual_Ek}), but neither factorizable nor
cofactorizable, and hence not invertible there.

\begin{rmk}[Indexing]
\label{rmk:indexing}
Our shift $n$ is the shift of the cotangent bundle being quantized: $\DMod^n\left(X\right)$
is $\EE_n$-monoidal and quantizes $T^*\left[n\right]X$.
In \cite{Ber:Sph} the same category is denoted
$\mathrm{Sph}\left(X , n-1\right)$, and is called the category of
\emph{$n$-shifted $\cD$-modules}.
\end{rmk}

Throughout this section, $k$ is algebraically closed, as in \cite{Ber:ShvCatH}. Following
\cite[Section 2]{Ber:Sph}, a \emph{stack} means a
derived algebraic stack over $k$ which is quasi-compact, locally
of finite presentation, with affine diagonal and perfect cotangent complex.
We also assume stacks are perfect in the sense of \cite{BZFN}, as required in \cite{HoLi};
this is preserved by the mapping stacks below \cite[Remark 3.1.1]{HoLi}.
A map of stacks $Y\to Z$ is \emph{bounded} if $S\times_Z Y$ is bounded for every bounded
affine scheme $S\to Z$ \cite[Section 2.1.5]{Ber:center}, \cite[Section 2.2.3]{HoLi}.

\subsection{Ind-coherent sheaves and formal completions}
\label{sec:IndCoh}

We make a brief aside to define the notion of ind-coherent sheaves, as it will be relevant
throughout this section.

Let $X$ be a derived scheme locally almost of finite type, and write
$\Coh\left(X\right) \subset \QC\left(X\right)$ for the complexes with bounded coherent
cohomology. The category of ind-coherent sheaves on $X$ is defined to be the
ind-completion of this:
\begin{equation*}
\IndCoh\left(X\right)\ceqq \Ind\left(\Coh\left(X\right)\right)
\end{equation*}
This is extended to stacks by taking the limit over affines along $!$-pullback
\cite[Part II]{GR:v1}.
There is a functor
$\Upsilon_X \colon \QC\left(X\right)\to \IndCoh\left(X\right)$,
$F \mapsto F\otimes \omega_X$, intertwining $f^*$ with $f^!$.
It is an equivalence when $X$ is smooth, and fully faithful when $X$ is bounded, i.e.\
eventually coconnective (since then $\omega_X$ is coherent with
$\End\left(\omega_X\right)\simeq\cO_X$). Quasi-smooth stacks are bounded, but taking loops
quickly produces unbounded stacks: $\cL$ sends smooth stacks to quasi-smooth ones, and
quasi-smooth but non-smooth stacks to unbounded ones \cite[Section 2.2]{Ber:Sph}.

Given a map of stacks $f\colon Y\to Z$, not necessarily a closed embedding,
the \emph{formal completion} of $Z$ along $f$ is
$Z^\wedge_Y \ceqq Z\times_{Z_\dR} Y_\dR$,
where $\left(-\right)_\dR$ is the de Rham stack \cite{GR:v2}.
For example, $\pt^\wedge_Y = Y_\dR$, and $Z^\wedge_Y\simeq Z$ whenever $f$ is an
isomorphism on underlying reduced classical stacks; e.g.\
$\left(\cL V\right)^\wedge_V\simeq \cL V$ for a vector space $V$.
Formal completions are not themselves stacks in the above sense. They only enter through
the arrow $Y\to Z$, and $\IndCoh_0$ below is really a sheaf theory on such arrows.

\begin{defn}[{\cite{AG:crystal}, \cite[Section 3]{Ber:center}}]
\label{defn:IndCoh0}
Let $Y\to Z$ be a map of stacks with $Y$ bounded. Then
$\IndCoh_0\left(Z^\wedge_Y\right)$ is defined as the full subcategory of
$\IndCoh\left(Z^\wedge_Y\right)$ given by the fiber product
\begin{equation*}
\IndCoh_0\left(Z^\wedge_Y\right) \ceqq
\IndCoh\left(Z^\wedge_Y\right)\times_{\IndCoh\left(Y\right)} \QC\left(Y\right) \ ,
\end{equation*}
where the maps are $!$-pullback and $\Upsilon_Y$.
If $Y$ is unbounded, this fiber product is not to be considered, and the definition is
modified as in \cite[Definition 4.1.6]{Ber:center}.
\end{defn}

The notation is abusive, as $\IndCoh_0\left(Z^\wedge_Y\right)$ depends on the map $Y\to Z$.
When $Y$ is smooth, $\IndCoh_0\left(Z^\wedge_Y\right) = \IndCoh\left(Z^\wedge_Y\right)$.
$\IndCoh_0$ is functorial for correspondences \cite[Theorem 3.3.3]{Ber:center}, satisfies
K\"unneth \cite[Section 3.4]{Ber:center}, and satisfies descent \cite[Proposition 4.4.1]{Ber:center}, \cite[Proposition 2.3.6]{HoLi};
these are the inputs to the excision theorem \cite[Theorem 3.3.1]{HoLi}.

For $Z = \pt$ we get the \emph{exotic} $\cD$-modules $\IndCoh_0\left(Y_\dR\right)$ of
\cite{Ber:center}. These agree with the usual ones $\IndCoh\left(Y_\dR\right)$ for $Y$
smooth, but unlike the usual ones they can see derived structure.
The only case we will need is the following, in which the derived structure is invisible
after all.

\begin{lem}
\label{lem:DMod0_qsmooth}
Let $Y$ be a bounded stack and $f\colon Y\to Y'$ a map to a smooth stack which
induces an isomorphism $Y_\dR\simeq Y'_\dR$. Then the forgetful functor is an equivalence
\begin{equation*}
\IndCoh_0\left(Y_\dR\right) \lto{\simeq} \IndCoh\left(Y_\dR\right)\simeq
\IndCoh\left(Y'_\dR\right) = \DMod\left(Y'\right) \ .
\end{equation*}
\end{lem}

\begin{proof}
The $!$-pullback $\IndCoh\left(Y_\dR\right)\to\IndCoh\left(Y\right)$ factors as
$\IndCoh\left(Y'_\dR\right)\to\IndCoh\left(Y'\right)\lto{f^!}\IndCoh\left(Y\right)$.
Since $Y'$ is smooth, $\IndCoh\left(Y'\right) = \Upsilon_{Y'}\QC\left(Y'\right)$, and
$f^!\Upsilon_{Y'}\simeq\Upsilon_Y f^*$.
So every object of $\IndCoh\left(Y_\dR\right)$ has underlying ind-coherent sheaf in the
image of $\Upsilon_Y$, which is fully faithful as $Y$ is bounded. Hence the
projection from the fiber product in Definition \ref{defn:IndCoh0} is an equivalence.
\end{proof}

\begin{exm}
\label{exm:DMod0_qsmooth}
Lemma \ref{lem:DMod0_qsmooth} applies to $Y = Y'\times U\left[-1\right]\to Y'$ for
$Y'$ smooth and $U$ a vector space, where $U\left[-1\right]\ceqq\pt\times_U\pt$; and to
$Y = \mathfrak{g}\left[-1\right]/G\to BG$ for an affine algebraic group $G$.
\end{exm}

For a finite-dimensional vector space $U$, write $U\left[1\right]\ceqq BU$ and
$U\left[-1\right]\ceqq\pt\times_U\pt$ for the corresponding \emph{odd vector spaces}. Their
functions are exterior algebras: For $\epsilon=\pm1$,
$\cO\left(U\left[\epsilon\right]\right)\simeq\Lambda^\bullet\left(U^\vee\left[-\epsilon\right]\right)$.

\begin{prop}[Odd vector spaces]
\label{prop:odd}
Let $U$ be a finite-dimensional vector space.
\begin{enumerate}
\item The Weyl algebra of $U\left[\epsilon\right]$, $\epsilon=\pm1$, i.e.\ the algebra generated by the linear
functions $U^\vee\left[-\epsilon\right]$ and the constant vector fields $U\left[\epsilon\right]$, subject to
the graded commutation relations $\left[\xi,x\right]=\langle\xi,x\rangle$, is the Clifford algebra
$\mathrm{Cl}$ of the hyperbolic pairing on $U^\vee\left[-\epsilon\right]\dsum U\left[\epsilon\right]$. Its
action on $\cO\left(U\left[\epsilon\right]\right)$ is an isomorphism
$\mathrm{Cl}\simeq\End_k\left(\Lambda^\bullet\left(U^\vee\left[-\epsilon\right]\right)\right)$, so
$\mathrm{Cl}$ is Morita trivial: $\Mod_{\mathrm{Cl}}\simeq\Vect$.
\item Correspondingly, for $Y$ smooth and $U_1,U_2$ finite-dimensional vector spaces,
\begin{equation*}
\DMod^0\left(Y\times U_1\left[1\right]\times U_2\left[-1\right]\right)\simeq\DMod\left(Y\right) \ .
\end{equation*}
In particular, $\DMod^0$ of an odd vector space is $\Vect$.
\end{enumerate}
\end{prop}

\begin{proof}
For (1): The generators are odd, so graded commutators are anticommutators, and the relations
are $x^2=0$, $\xi^2=0$ and $\xi x+x\xi=\langle\xi,x\rangle$. These are the relations of the Clifford
algebra of the hyperbolic pairing, acting on $\Lambda^\bullet\left(U^\vee\left[-\epsilon\right]\right)$
by multiplication and contraction. This is the Fock representation, which is faithful, and both
sides have dimension $2^{2\dim U}$, so it is an isomorphism. Since
$\Lambda^\bullet\left(U^\vee\left[-\epsilon\right]\right)$ is a non-zero perfect complex, it is a
compact generator of $\Vect$, and $\Mod_{\End\left(P\right)}\simeq\Vect$ for $P$ such a generator.

For (2): $Y\times U_1\left[1\right]$ is smooth, so Lemma \ref{lem:DMod0_qsmooth} (Example
\ref{exm:DMod0_qsmooth}) gives
\begin{equation*}
\DMod^0\left(Y\times U_1\left[1\right]\times U_2\left[-1\right]\right)\simeq
\DMod\left(Y\times U_1\left[1\right]\right) \ .
\end{equation*}
It remains to show that pullback along the projection
$Y\times BU_1\to Y$ is an equivalence. By smooth descent for $\cD$-modules \cite{GR:v2} along
$Y\to Y\times BU_1$, whose \v{C}ech nerve is $Y\times U_1^{\bullet}$, the category
$\DMod\left(Y\times BU_1\right)$ is the limit of the cosimplicial category
$\DMod\left(Y\times U_1^\bullet\right)$. All face maps of this \v{C}ech nerve commute with the
projections to $Y$, and pullback along each projection $Y\times U_1^p\to Y$ is fully faithful, by
the projection formula and $H^*_\dR\left(\AA^m\right)\simeq k$. The component in degree $p$ of an
object of the limit is a pullback of its component in degree $0$, so it lies in the essential
image of $\DMod\left(Y\right)$, and on these essential images the cosimplicial category is
equivalent to the constant one with value $\DMod\left(Y\right)$. Since $\Delta$ is weakly
contractible, the limit is $\DMod\left(Y\right)$.
\end{proof}

\subsection{Shifted \texorpdfstring{$\cD$}{D}-modules}

\begin{defn}[\cite{Ber:Sph}]
\label{defn:DMod_n}
Given a stack $X$ and $n\geq 0$,
the category of $n$-shifted $\cD$-modules on $X$ is
\begin{equation*}
\DMod^n\left(X\right) \ceqq
\IndCoh_0\left(
\left(
X^{S^{n-1}}
\right)^{\wedge}_X
\right) \ ,
\end{equation*}
where $X\to X^{S^{n-1}}$ is the inclusion of constant maps, and $S^{-1}\ceqq\emptyset$.
\end{defn}

The $\EE_{n}$-monoidal structure comes from the fact that this is (sheaves on the
completed) mapping stack from $S^{n-1}$: The multiplication is $\IndCoh_0$-pull-push
along the restriction maps from $X^{P_n}$, where $P_n\ceqq\left(D^n\right)^{\circ\circ}$
is the $n$-dimensional pair of pants, and similarly for more little disks.
For a discussion of the $\mathbb{E}_n$-monoidal structure, we recommend
\cite{Toen:branes}.
$SO\left(n\right)$ acts on all of this data, so $\DMod^n\left(X\right)$ defines an
\emph{oriented} $n$-dimensional TQFT, computed by factorization homology, via the cobordism
hypothesis \cite{L:Cob}, \cite[Section 3.2]{Ber:Sph}. For $n=2$, the only $3$-dimensional theory we
use is the framed one (Lemma \ref{lem:inv_surfaces}).
In the notation of \cite{HoLi}, $\DMod^n\left(X\right)=\mathrm{Hecke}_n\left(X,\pt\right)$.
There the $\EE_n$-structure is constructed as the restriction to disks of a symmetric monoidal
functor out of all $n$-manifolds and open embeddings \cite[Sections 3.1--3.2]{HoLi}, so no
framing is needed.

\begin{exm}
\label{exm:DMod_small_n}
\begin{enumerate}
\item For $n = 0$ we obtain the exotic $\cD$-modules $\DMod^0\left(X\right) =
\IndCoh_0\left(X_\dR\right)$, which are the usual ones when $X$ is smooth.
\item For $n = 1$ we obtain the monoidal category
\begin{equation*}
\DMod^1\left(X\right) = \IndCoh_0\left(\left(X\times X\right)^\wedge_X\right)
\end{equation*}
of \cite{Ber:center}.
\item For $n=2$ we obtain the $\EE_2$-monoidal category
\begin{equation*}
\DMod^2\left(X\right) = \IndCoh_0\left(\left(\cL X\right)^\wedge_X\right).
\end{equation*}
For a vector space $V$ this is $\IndCoh\left(\cL V\right)
= \IndCoh\left(V\times V\left[-1\right]\right)$, and for $X = BG$ it is
$\IndCoh\left(\widehat{G}_e/G\right)$, where $G$ acts on the formal group $\widehat G_e$ by
conjugation.
\end{enumerate}
\end{exm}

\begin{rmk}
\label{rmk:DMod_center}
\begin{enumerate}
\item For $X$ a smooth scheme, $X\to\cL X$ is a nil-isomorphism, so
$\DMod^2\left(X\right)\simeq\IndCoh\left(\cL X\right)$. The convolution category
$\QC\left(\cL X\right)$ of Section \ref{sec:QC} is obtained from it by requiring the singular
support to vanish \cite{AG:singsupp}, i.e.\ by completing along the zero section of
$T^*\left[2\right]X$.
\item Both are Drinfeld centers. $\QC\left(\cL X\right)$ is the center of $\QC\left(X\right)$
(Theorem \ref{thm:BZFN_loops}), i.e.\ of the $2$-category of sheaves of categories on $X$ when
$X$ is $1$-affine. Similarly, $\DMod^2\left(X\right)\simeq\IndCoh\left(\cL X\right)$ is the center of the
$2$-category $2\IndCoh\left(X\right)$ of \cite[Section 1]{Ste:IndCoh}, and it is the value on $S^1$ of
the three-dimensional TQFT built from $2\IndCoh\left(X\right)$ \cite[Corollary 1.3]{Ste:IndCoh}, a
version of Rozansky--Witten theory with target $T^*\left[2\right]X$ (cf.\ \cite[Section D.6]{BZSV}).
\item The symmetric monoidal category $\DMod\left(X\right)$ of ordinary $\cD$-modules acts
centrally on $\DMod^n\left(X\right)$ for every $n\geq1$, so $\DMod^n\left(X\right)$ sheafifies over
$X_\dR$. So invertibility can only be expected when $X$ has no de Rham cohomology, in the sense
that $\DMod\left(X\right)\simeq\Vect$. Proposition \ref{prop:BG}, where
$\cZ_2\left(\DMod^2\left(BG\right)\right)\simeq\DMod\left(BG\right)$, makes this precise for $X=BG$,
and Theorem \ref{thm:BW} shows that invertibility can fail even when $\DMod\left(X\right)\simeq\Vect$.
\end{enumerate}
\end{rmk}

\subsubsection{Factorization homology}

The factorization homology of shifted $\cD$-modules is computed by an excision theorem for
$\IndCoh_0$ \cite[Theorem 3.3.1]{HoLi}, cf.\ \cite[Claim 3.1]{Ber:Sph}.

\begin{prop}[{\cite[Theorem 3.3.1 and Corollary 1.3.12]{HoLi}}]
\label{prop:int_DMod}
Let $M$ be a closed $m$-manifold, $m\leq n$. Then, integrating over $M\times\RR^{n-m}$,
\begin{equation*}
\int_M\DMod^n\left(X\right)\simeq\DMod^{n-m}\left(X^M\right) \ .
\end{equation*}
In particular, $\int_\Sigma\DMod^2\left(X\right)\simeq\DMod^0\left(X^\Sigma\right)=
\IndCoh_0\left(\left(X^\Sigma\right)_\dR\right)$ for a closed oriented surface $\Sigma$.
\end{prop}

This is \cite[Corollary 1.3.12]{HoLi} with $Z=\pt$, since
$\partial\left(M\times D^{n-m}\right)=M\times S^{n-m-1}$ and
$X^{M\times S^{n-m-1}}\simeq\left(X^M\right)^{S^{n-m-1}}$; cf.\ also \cite[Claim 3.1 and
Theorem 5.3]{Ber:Sph}. For $M$ oriented, it is the quantization of the AKSZ
statement \cite{PTVV}
\begin{equation*}
\Maps\left(M , T^*\left[n\right]X\right) \simeq T^*\left[n-m\right]X^M \ .
\end{equation*}

\subsection{Rigidity and \texorpdfstring{$3$}{3}-dualizability}
\label{sec:rigid}

For $n\geq 1$, gluing $S^{n-1}$ from two disks along $S^{n-2}$ gives
$X^{S^{n-1}}\simeq X\times_{X^{S^{n-2}}}X$, so
\begin{equation*}
\DMod^n\left(X\right)\simeq\IndCoh_0\left(\left(X\times_{X^{S^{n-2}}}X\right)^\wedge_X\right) \ ,
\end{equation*}
and the monoidal structure coming from any one of the $n$ directions (these are all equivalent
by $SO\left(n\right)$-invariance) is convolution relative to $X^{S^{n-2}}$ \cite{Ber:center}.
This is the $\EE_1$-Hecke category $\mathrm{Hecke}_1\left(X,X^{S^{n-2}}\right)$ of \cite{HoLi}.
Recall that a monoidal category in $\Pre$ is \emph{rigid} if it is compactly generated, its
unit is compact, and every compact object has a left and a right dual
\cite[Chapter 1, Section 9]{GR:v1}, \cite{G:ShvCat}.
For $n=1$, the following is observed in \cite{Ber:center}; the same argument applies relative
to $X^{S^{n-2}}$.

\begin{prop}
\label{prop:rigid}
Let $X$ be a bounded stack and $n\geq1$. Then
$\DMod^n\left(X\right)$ is rigid, and it is compactly generated by the objects
$\iota\left(P\right)$ for $P\in\Perf\left(X\right)$, where
$\iota\colon\QC\left(X\right)\to\DMod^n\left(X\right)$ is pushforward along the diagonal.
\end{prop}

\begin{proof}
Write $Z\ceqq X^{S^{n-2}}$ and $\Delta\colon X\to\left(X\times_Z X\right)^\wedge_X$ for the diagonal. It induces an
equivalence on de Rham stacks, so it is an inf-schematic nil-isomorphism; in particular it is
ind-proper, so $\Delta_*\dashv\Delta^!$, and $\Delta^!$ is conservative \cite[Part~I]{GR:v2}.
\begin{enumerate}
\item By our conventions $X$ is QCA, so in characteristic zero $\QC\left(X\right)$ is
compactly generated by $\Perf\left(X\right)$, and $\cO_X$ is compact \cite{DG:finiteness}.
\item The functor
\begin{equation*}
\iota = \Delta_*\circ\Upsilon_X
\end{equation*}
lands in $\IndCoh_0$, and is monoidal from $\left(\QC\left(X\right),\otimes\right)$ to the
convolution structure \cite{Ber:center}.
Since $\Upsilon_X$ is fully faithful, $\IndCoh_0$ is the full subcategory of
\begin{equation*}
\IndCoh\left(\left(X\times_ZX\right)^\wedge_X\right)
\end{equation*}
on objects $F$ with
$\Delta^!F\in\Upsilon_X\QC\left(X\right)$. So the right adjoint of $\iota$ is
$\Upsilon_X^{-1}\circ\Delta^!$, the projection in Definition \ref{defn:IndCoh0}, which is
continuous and conservative.
\item Hence $\iota$ preserves compact objects, and the $\iota\left(P\right)$ for
$P\in\Perf\left(X\right)$ form a set of compact generators. In particular the unit
$\iota\left(\cO_X\right)$ is compact.
\item Each $\iota\left(P\right)$ has left and right duals, since $P$ does and $\iota$ is
monoidal. Objects with left (resp.\ right) duals form a thick subcategory, and the compact
objects form the thick subcategory generated by the $\iota\left(P\right)$. So every compact
object has duals. \qedhere
\end{enumerate}
\end{proof}

\begin{cor}
\label{cor:3dual}
Let $X$ be bounded and $\cA = \DMod^2\left(X\right)$. Then $\cA$ and
$\cA^e\simeq\DMod^2\left(X^2\right)$ are rigid.
If moreover $\cL X$ is bounded, e.g.\ if $X$ is smooth, then
$\HC\left(\cA\right)\simeq\DMod^1\left(\cL X\right)$ is rigid,
$\cA$ is dualizable in $\Pre$ and as a module over $\cA^e$ and over $\HC\left(\cA\right)$, and $\cA$
is $3$-dualizable in $\umor_2\left(\Pre\right)$.
\end{cor}

\begin{proof}
Apply Proposition \ref{prop:rigid} to $\DMod^2\left(X\right)$, $\DMod^2\left(X^2\right)$ and
$\DMod^1\left(\cL X\right)$. The identification of
$\HC\left(\cA\right)$ as a monoidal category is \cite[Corollary 5.2]{Ber:Sph}; it is also
\cite[Theorem 3.3.1]{HoLi} for $Z=\pt$ and $M=S^1\times\RR$, with stacking of cylinders
inducing convolution (cf.\ \cite[Example 1.3.8]{HoLi}).
Since $\cA$ is compactly generated, it is dualizable in $\Pre$. A module over a rigid
monoidal category which is dualizable in $\Pre$ is dualizable as a module
\cite[Proposition 2.1.6]{Ber:ShvCatH}, cf.\ \cite[Chapter 1, Section 9]{GR:v1}. Now apply Theorem \ref{thm:inv_E2}.
\end{proof}

\subsection{Higher centers and surfaces}
\label{sec:DMod_centers}

Let $\cA$ be an $\EE_2$-algebra in $\Pre$ which is $SO\left(2\right)$-invariant, so that
factorization homology over oriented surfaces makes sense.

\begin{lem}
\label{lem:inv_surfaces}
If $\cA$ is invertible in $\umor_2\left(\Pre\right)$, then $\int_{T^2}\cA$ is an invertible
object of $\Pre$, for any framing of $T^2$.
\end{lem}

\begin{proof}
The framed TQFT defined by $\cA$ has invertible value on the point, so it factors through the
maximal Picard $\infty$-groupoid of the target, and it takes invertible values on all framed
bordisms.
Its value on the framed torus is $\int_{T^2}\cA$ \cite{Sch:thesis}, an object of
$\Omega^2\umor_2\left(\Pre\right)\simeq\Pre$ (Section \ref{sec:mor}).
\end{proof}

This is a condition on the values $\int_\Sigma\cA$, \emph{not} on $\cA$ itself: An invertible
$\EE_2$-algebra need not be invertible in $\Pre$. For example, a modular tensor category is
invertible in $\umor_2\left(\Pre\right)$ \cite{BJSS} without being equivalent to $\Vect$, but its
factorization homology over every closed surface is $\Vect$ \cite[Theorem 9.7]{KT:4d}
(both in the abelian setting, cf.\ Remark \ref{rmk:abelian}).

\begin{lem}
\label{lem:unit_Vect}
Let $\cZ$ be a monoidal category in $\Pre$ which is equivalent to $\Vect$ as a plain
category. Then the unit map $\Vect\to\cZ$ is an equivalence.
\end{lem}

\begin{proof}
Identify $\cZ$ with $\Vect$. A colimit-preserving $k$-bilinear functor is of the form
\begin{equation*}
\Vect\times\Vect\to\Vect\ ,\quad \left(x,y\right)\mapsto x\otimes y\otimes E
\end{equation*}
with
$E\ceqq k\otimes_\cZ k$. If $\one$ is the unit, then $\one\otimes_\cZ k\simeq k$ says
$\one\otimes E\simeq k$, so $\one$ and $E$ are invertible in $\Vect$, i.e.\ shifts of $k$.
So $x\mapsto x\otimes\one$ is an equivalence.
\end{proof}

The conditions of Definition \ref{defn:inv_E2} concern the higher centers
\begin{equation*}
\cZ_j\left(\cA\right)=\End_{\int_{S^{j-1}\times\RR^{3-j}}\cA}\left(\cA\right)
\end{equation*}
rather than
factorization homology. For shifted $\cD$-modules, $\cZ_j$ is factorization homology over the
$j$-sphere.

\begin{lem}
\label{lem:H_duals}
\leavevmode
\begin{enumerate}
\item For $X$ bounded and $n\geq1$, the monoidal category $\cC=\DMod^n\left(X\right)$ is its own
right dual as a left $\cC^e$-module.
\item Let $f\colon\mathcal{Y}\to\mathcal{Z}$ be a bounded map of bounded stacks. Then the
transfer bimodules between $\DMod^1\left(\mathcal{Y}\right)$ and $\DMod^1\left(\mathcal{Z}\right)$
\begin{equation*}
\mathbb{T}_{\mathcal{Y}\to\mathcal{Z}}\ceqq
\IndCoh_0\left(\left(\mathcal{Y}\times\mathcal{Z}\right)^\wedge_{\mathcal{Y}}\right)
\quad\text{and}\quad
\mathbb{T}_{\mathcal{Z}\leftarrow\mathcal{Y}}\ceqq
\IndCoh_0\left(\left(\mathcal{Z}\times\mathcal{Y}\right)^\wedge_{\mathcal{Y}}\right)
\end{equation*}
are left and right dual to each other \cite{Ber:ShvCatH}.
\end{enumerate}
\end{lem}

\begin{proof}
(1) The category $\cC$ is rigid by Proposition \ref{prop:rigid}. It is also pivotal, since the
proof of \cite[Lemma 3.5.4]{Ber:center} only uses rigidity and the self-duality of
\cite[Proposition 3.2.4]{Ber:center}. So $\cC$ is very rigid \cite[Section 5.1.4]{Ber:center},
i.e.\ the self-duality $\cC\simeq\cC^\vee$ is $\left(\cC,\cC\right)$-bilinear.
By \cite[Corollary 2.1.7]{Ber:ShvCatH}, the right dual of $\cC$ as a module over the rigid
category $\cC^e$ is $\cC^\vee$, so it is $\cC$.

(2) These are the values of the $2$-functor of \cite[Theorem 6.5.3]{Ber:ShvCatH} on the correspondences
$\left[\mathcal{Y}=\mathcal{Y}\to\mathcal{Z}\right]$ and $\left[\mathcal{Z}\leftarrow\mathcal{Y}=\mathcal{Y}\right]$.
Tensoring with them realizes pullback $f^*$ and pushforward $f_*$ between
$\DMod^1\left(\mathcal{Z}\right)$- and $\DMod^1\left(\mathcal{Y}\right)$-modules
\cite[Corollary 6.5.4]{Ber:ShvCatH}, and $f_!\dashv f^*\dashv f_*$ with
$f_!\simeq f_*$ \cite[Corollary 6.5.5]{Ber:ShvCatH}.
\end{proof}

\begin{theorem}
\label{thm:centers}
Let $X$ be smooth and $\cA = \DMod^2\left(X\right)$. For $j=0,1$, and for $j=2$ if the constant
loops $X\to\cL X$ are bounded,\footnote{This holds e.g.\ for $X=BG$ with $G$ an affine algebraic
group, since the pullback of $BG\to G/G$ along the smooth cover $G\to G/G$ is $e\to G$. It fails
for smooth schemes of positive dimension, for which $X^{S^2}\simeq T\left[-2\right]X$ is unbounded
(cf.\ \cite[Section 1.3.13]{HoLi}). We expect this assumption to be unnecessary.}
there are equivalences in $\Pre$
\begin{equation*}
\cZ_j\left(\cA\right)\simeq\int_{S^j\times\RR^{2-j}}\cA\simeq\DMod^{2-j}\left(X^{S^j}\right) \ .
\end{equation*}
For $j=0$ the middle term is $\cA\otimes\cA$, and for $j=1$ it is $\HC\left(\cA\right)$.
\end{theorem}

\begin{proof}
The second equivalence is Proposition \ref{prop:int_DMod}. For the first equivalence, let
$\cB_j\ceqq\int_{S^{j-1}\times\RR^{3-j}}\cA$, so that $\cZ_j\left(\cA\right)=\End_{\cB_j}\left(\cA\right)$.
If $\cA$ has a right dual $\cA^R$ as a left $\cB_j$-module, then
\begin{equation*}
\End_{\cB_j}\left(\cA\right)\simeq\cA^R\otimes_{\cB_j}\cA \ ,
\end{equation*}
so it suffices to find $\cA^R$.

For $j=0$, $\cA^R=\cA^\vee\simeq\cA$ \cite[Proposition 3.2.4]{Ber:center}, so
$\cZ_0\left(\cA\right)\simeq\cA\otimes\cA$.

For $j=1$, $\cB_1=\cA^e$, so $\cA^R\simeq\cA$ by Lemma
\ref{lem:H_duals}(1), and $\cZ_1\left(\cA\right)\simeq\cA\otimes_{\cA^e}\cA=\HC\left(\cA\right)$.

For $j=2$, $\cB_2=\HC\left(\cA\right)\simeq\DMod^1\left(\cL X\right)$ (Corollary \ref{cor:3dual}).
Gluing the collar $S^1\times\RR$ onto the disk gives the correspondence
$\cL X\leftarrow X\to\pt$ of mapping stacks \cite[Section 3.1]{HoLi}, so $\cA$ is the value of
the $2$-functor of \cite[Theorem 6.5.3]{Ber:ShvCatH} on it. Factoring through $X$,
\begin{equation*}
\cA\simeq\mathbb{T}_{\cL X\leftarrow X}\otimes_{\DMod^1\left(X\right)}\mathbb{T}_{X\to\pt} \ .
\end{equation*}
By Lemma \ref{lem:H_duals}(2), applied to $X\to\cL X$ and $X\to\pt$, the right dual $\cA^R$ is the
value on the reversed correspondence $\pt\leftarrow X\to\cL X$. So $\cZ_2\left(\cA\right)$ is the
value on the composite $\pt\leftarrow X\times_{\cL X}X\to\pt$, which is
$\DMod^0\left(X^{S^2}\right)$.
\end{proof}

\begin{rmk}
\label{rmk:centers_spheres}
By Theorem \ref{thm:centers}, for $X$ as there the comparison maps of Corollary
\ref{cor:SSS_inv} for $\cA=\DMod^2\left(X\right)$ are, up to the equivalences of that theorem
(which are not claimed to be compatible with these maps), maps between factorization homologies
over spheres,
\begin{equation*}
\int_{S^{i-1}\times\RR^{3-i}}\cA\longrightarrow\int_{S^{2-i}\times\RR^{i}}\cA \ .
\end{equation*}
For $i=0$ this is $\Vect\to\int_{S^2}\cA$ (non-degeneracy), for $i=1$ it is
$\cA\otimes\cA^{\sigma\rev}\simeq\int_{S^0\times\RR^2}\cA\to\int_{S^1\times\RR}\cA$
(factorizability), and for $i=2$ it is $\int_{S^1\times\RR}\cA\to\int_{S^0\times\RR^2}\cA$
(cofactorizability). So factorizability and cofactorizability compare the same two categories,
$\DMod^2\left(X^2\right)$ and $\DMod^1\left(\cL X\right)$, in opposite directions.
\end{rmk}

\subsection{The critical case: \texorpdfstring{$X = V\left[1\right]$}{X = V[1]}}

Let $V$ be a finite-dimensional vector space, and set $X\ceqq V\left[1\right]$ and
$\cA \ceqq \DMod^2\left(X\right)$.

\begin{theorem}
\label{thm:BW}
For a closed oriented surface $\Sigma_g$ of genus $g$,
\begin{equation*}
\int_{\Sigma_g}\DMod^2\left(V\left[1\right]\right)\simeq\DMod\left(V^{2g}\right) \ .
\end{equation*}
Consequently:
\begin{enumerate}
\item $\cA$ is $3$-dualizable and non-degenerate: $\cZ_2\left(\cA\right)\simeq\Vect$;
\item if $V\neq 0$ then $\int_{T^2}\cA$ is not invertible, so $\cA$ is not invertible;
\item if $V\neq 0$ then $\cA$ is neither factorizable nor cofactorizable.
\end{enumerate}
\end{theorem}

\begin{proof}
Since $V\left[1\right]$ is a linear stack, we have
$X^{\Sigma_g}\simeq V\otimes C^*\left(\Sigma_g\right)\left[1\right]$, i.e.
\begin{equation*}
\left(V\left[1\right]\right)^{\Sigma_g}\simeq
V\left[1\right]\dsum V^{2g}\dsum V\left[-1\right]
\simeq V\left[1\right]\times V^{2g}\times V\left[-1\right] \ .
\end{equation*}
This is quasi-smooth, hence bounded, and the projection to the smooth stack $V\left[1\right]\times V^{2g}$ is an
isomorphism on de Rham stacks. The other stacks involved, $V\left[1\right]$,
$\cL\left(V\left[1\right]\right)\simeq V\times V\left[1\right]$ and
$V\left[1\right]\times V\left[1\right]$, are smooth. So every stack that appears is bounded, and Definition
\ref{defn:IndCoh0} applies as stated. (In general, this is the Hecke pair condition of
\cite[Definition 3.4.2]{HoLi}, which holds for $BH\to BG$ for any homomorphism of affine
algebraic groups \cite[Lemma 3.4.4]{HoLi}, here $BV\to\pt$.) By Proposition \ref{prop:int_DMod} and Proposition \ref{prop:odd}(2), with $Y=V^{2g}$ and
$U_1=U_2=V$,
\begin{equation*}
\int_{\Sigma_g}\cA\simeq\DMod^0\left(\left(V\left[1\right]\right)^{\Sigma_g}\right)
\simeq\DMod\left(V^{2g}\right) \ .
\end{equation*}

For (1): $V\left[1\right]$ is smooth, so $\cA$ is $3$-dualizable by Corollary \ref{cor:3dual}.
Since $V\left[1\right]=BV$, the constant loops $V\left[1\right]\to\cL\left(V\left[1\right]\right)$
are bounded (see the footnote to Theorem \ref{thm:centers}), so
$\cZ_2\left(\cA\right)\simeq\int_{S^2}\cA\simeq\Vect$ by Theorem \ref{thm:centers} and the
case $g=0$. The unit map is then an equivalence by Lemma \ref{lem:unit_Vect}.
For (2) take $g=1$: The value of the TQFT on the torus is
$\int_{T^2}\cA\simeq\DMod\left(V^2\right)\simeq\Mod_D$, with $D$ the Weyl algebra
of $V^2$. This object of $\Pre$ is not invertible: Otherwise
$\Mod_{D^\op\otimes D}\simeq\Mod_D^\vee\otimes\Mod_D\simeq\Vect$, but the compact object
$D^\op\otimes D$ has infinite-dimensional endomorphisms. By Lemma
\ref{lem:inv_surfaces}, $\cA$ is not invertible.

For (3), call a compactly generated category \emph{proper} if all mapping complexes between
compact objects have finite-dimensional total cohomology. This is invariant under equivalence in
$\Pre$, and it holds if there is a compact generator with finite-dimensional endomorphisms. By
Remark \ref{rmk:centers_spheres}, it suffices to show that $\int_{S^0\times\RR^2}\cA$ is proper
and $\int_{S^1\times\RR}\cA$ is not.
\begin{itemize}
\item $\int_{S^0\times\RR^2}\cA\simeq\DMod^2\left(X^2\right)$ is proper. Write $U\ceqq V^2$, so
$X^2\simeq BU$ and $\left(\cL BU\right)^\wedge_{BU}\simeq\widehat U_0\times BU$, with $\widehat U_0$
the formal completion of $U$ at $0$ (Example \ref{exm:DMod_small_n}), so
$\DMod^2\left(X^2\right)\simeq\IndCoh\left(\widehat U_0\right)\otimes\QC\left(BU\right)$. Since $U$
is unipotent, $\Perf\left(BU\right)$ is generated by $\cO_{BU}$, so by Proposition
\ref{prop:rigid} this is compactly generated by
$\iota\left(\cO_{BU}\right)=k_0\boxtimes\cO_{BU}$. Its endomorphisms are
$\End\left(k_0\right)\otimes C^*\left(BU\right)\simeq
\Lambda^\bullet\left(U\left[-1\right]\right)\otimes\Lambda^\bullet\left(U^\vee\left[-1\right]\right)$,
which is finite-dimensional.
\item $\int_{S^1\times\RR}\cA\simeq\DMod^1\left(\cL X\right)$ is
not proper. Its unit $\iota\left(\cO_{\cL X}\right)$ is compact by Proposition \ref{prop:rigid}.
The action of $\DMod^1\left(\cL X\right)$ on $\IndCoh\left(\cL X\right)$ by integral
transforms restricts along $\iota$ to the action of $\QC\left(\cL X\right)$ by tensoring, so the
composite
$\Gamma\left(\cL X,\cO\right)\to\End\left(\iota\left(\cO_{\cL X}\right)\right)\to
\End\left(\omega_{\cL X}\right)\simeq\Gamma\left(\cL X,\cO\right)$ is the identity. Hence
$\Gamma\left(\cL X,\cO\right)\simeq\cO\left(V\right)\otimes\Lambda^\bullet\left(V^\vee\left[-1\right]\right)$,
which is infinite-dimensional as $V\neq0$, is a retract of the endomorphisms of a compact object.
\end{itemize}
So neither the factorizability map
$\cA\otimes\cA^{\sigma\rev}\simeq\int_{S^0\times\RR^2}\cA\to\cZ_1\left(\cA\right)\simeq\int_{S^1\times\RR}\cA$
nor the cofactorizability map
$\HC\left(\cA\right)=\int_{S^1\times\RR}\cA\to\cZ_0\left(\cA\right)\simeq\int_{S^0\times\RR^2}\cA$
is an equivalence.
\end{proof}

Theorem \ref{thm:BW} is an instance of a parity phenomenon. For a linear
stack $X=U$, i.e.\ a graded vector space, $X^M\simeq U\otimes H^*\left(M\right)$. By
Proposition \ref{prop:odd}, odd directions contribute Clifford algebras, which are Morita trivial,
while even directions contribute Weyl algebras, which are not. Since $H^*\left(S^0\right)$ and
$H^*\left(S^2\right)$ are concentrated in even degrees, $U\otimes H^*\left(S^0\right)$ and
$U\otimes H^*\left(S^2\right)$ are odd when $U$ is, so the $S^2$ test of non-degeneracy is passed.
But $H^*\left(S^1\right)$ and $H^*\left(T^2\right)$ have odd classes, so $U\otimes H^*\left(S^1\right)$
and $U\otimes H^*\left(T^2\right)$ have non-zero even part unless $U=0$. For $U=V\left[1\right]$,
this even part is the factor $V$ of $\cL X\simeq V\times V\left[1\right]$, which makes
$\int_{S^1\times\RR}\cA$ non-proper and so breaks both factorizability and cofactorizability, and
the factor $V^{2g}$ of
$\left(V\left[1\right]\right)^{\Sigma_g}\simeq V\left[1\right]\times V^{2g}\times V\left[-1\right]$,
which makes the $T^2$ test fail.

\begin{prop}
\label{prop:BG}
Let $G$ be an affine algebraic group. Then
$\cZ_2\left(\DMod^2\left(BG\right)\right)\simeq\DMod\left(BG\right)$. Hence
$\DMod^2\left(BG\right)$ is non-degenerate if and only if $G$ is unipotent.
\end{prop}

\begin{proof}
Since $BG$ is smooth and $BG\to\cL BG\simeq G/G$ is bounded, Theorem \ref{thm:centers} applies. With
$BG^{S^2}\simeq\mathfrak{g}\left[-1\right]/G$ and Example \ref{exm:DMod0_qsmooth}, it gives
\begin{equation*}
\cZ_2\left(\DMod^2\left(BG\right)\right)\simeq\DMod^0\left(\mathfrak{g}\left[-1\right]/G\right)
\simeq\DMod\left(BG\right) \ .
\end{equation*}
By Lemma \ref{lem:unit_Vect}, it follows that $\DMod^2\left(BG\right)$ is non-degenerate
exactly when $\DMod\left(BG\right)\simeq\Vect$.
If $G$ is unipotent, then $\DMod\left(BG\right)\simeq\Vect$ as in the proof of Proposition
\ref{prop:odd}(2), since $G\simeq\AA^{\dim G}$ as a variety.
Conversely, suppose $\DMod\left(BG\right)\simeq\Vect$.
We may assume $k=\CC$: $G$ is defined over a finitely generated subfield, which embeds into $\CC$,
unipotence is invariant under field extension, and so is the condition that the unit
$\Vect\to\DMod\left(BG\right)$ is an equivalence, since $\DMod$ commutes with extension of scalars and
extension of scalars is conservative.
Let $K$ be a maximal compact subgroup of $G\left(\CC\right)$. The endomorphisms of the unit are
the de Rham cohomology of $BG$, which by the comparison theorem is
$H^*\left(BG\left(\CC\right);\CC\right)\simeq H^*\left(BK;\CC\right)$, since $G\left(\CC\right)$
deformation retracts onto $K$. If $K^\circ\neq1$, a $K$-invariant inner product on its Lie
algebra gives a non-zero class in $H^4\left(BK;\CC\right)$, so $K$ is finite. Then the identity
component $G^\circ$ is unipotent and
$\Gamma\ceqq\pi_0\left(G\right)\simeq\pi_0\left(G\left(\CC\right)\right)\simeq K$. Since
$\DMod\left(BG^\circ\right)\simeq\Vect$, pullback along $q\colon BG\to B\Gamma$ is fully faithful: By
the projection formula and base change along $\pt\to B\Gamma$, the unit $\Id\to q_*q^*$ is
tensoring with the de Rham cohomology of the fiber $BG^\circ$, which is $k$. So
$\Rep\left(\Gamma\right)=\DMod\left(B\Gamma\right)$ embeds fully faithfully into $\Vect$. In
$\Rep\left(\Gamma\right)$ distinct irreducibles have zero mapping complexes, whereas any two
non-zero objects of $\Vect$ have a non-zero mapping complex, so $\Gamma$ is trivial and $G$ is
unipotent.
\end{proof}

For example, $\DMod\left(B\Gamma\right)=\Rep\left(\Gamma\right)$ for $\Gamma$ finite,
$\DMod\left(BT\right)$ is modules over an exterior algebra for a torus $T$, and in general
$\DMod\left(BG\right)\simeq\DMod\left(BG^{\mathrm{red}}\right)$ for $G^{\mathrm{red}}\ceqq G/R_u\left(G\right)$
the reductive quotient.

\begin{cor}
\label{cor:BG_not_inv}
Let $G$ be a non-trivial affine algebraic group which is either not unipotent or abelian.
Then $\DMod^2\left(BG\right)$ is not invertible.
\end{cor}

\begin{proof}
If $G$ is not unipotent, then $\DMod^2\left(BG\right)$ is degenerate by Proposition
\ref{prop:BG}, so not invertible by Theorem \ref{thm:inv_E2}. If $G$ is abelian unipotent, then
$G$ is the additive group of a vector space $V$, so $BG\simeq V\left[1\right]$, and this is Theorem \ref{thm:BW}(2).
\end{proof}
